\chapter{Conclusions and Future Work}
\label{ch:future_work}

The work presented in this thesis establishes a mathematically sound,
performant, and truly language-agnostic framework for architectural
dependency extraction. By abstracting the extraction logic through a
heuristic-based approach on Concrete Syntax Trees (CST) and
formalizing a flexible Intermediate Representation
($\mathcal{D}_{IR}$), the system successfully overcomes the
scalability and configuration bottlenecks of previous iterations. The
validation phase demonstrated the robustness of the model across
different object-oriented and procedural languages (such as Java,
Python, C++, and Rust), properly resolving complex nested scopes,
generic types, and cross-module imports.

However, as the primary focus of this research was the
conceptualization, formalization, and validation of the theoretical
model, several avenues remain open for future improvements and
extensions. The following sections outline the most promising
directions for future research.

\section{Evolution of the Extraction Phase: srcML Integration}

While the current heuristic-based extraction directly from
Tree-sitter CSTs eliminates the need for manual dispatch tables, it
remains the most delicate phase of the pipeline. The framework relies
on structural heuristics and string-matching of node names (e.g.,
identifying nodes containing \texttt{struct} or \texttt{class}).
Consequently, it is somewhat fragile to unexpected grammar updates or
drastic structural variations in Tree-sitter parsers.

A highly promising future development involves replacing or
augmenting the first extraction phase with an XML-based parser such
as \textbf{srcML}. As discussed in the Related Work chapter, srcML
standardizes source code into an XML format while preserving all
original text and structural context. Transitioning to srcML would
completely decouple the extraction engine from Tree-sitter's AST
volatility, transforming the heuristic identification of
architectural entities into a purely declarative and robust process
(e.g., using standardized XPath queries to reliably extract
\texttt{<class>} or \texttt{<function>} nodes).

\section{Performance Optimization and Concurrency}

Throughout the development of this thesis, the primary objective was
ensuring the correctness and scalability of the mathematical model,
the resolution strategies, and the algebraic query engine.
Consequently, there has been limited focus on low-level performance
optimization, and the current prototype operates almost entirely sequentially.

Future iterations of the framework could achieve massive performance
gains by introducing concurrency. The Extraction Phase (Phase 1) is
intrinsically parallelizable, as each source file is parsed and
analyzed independently. Implementing a thread pool architecture to
process files concurrently would drastically reduce the initial
parsing time. Similarly, the Resolution Phase (Phase 2) can be
optimized by parallelizing independent query executions across
different namespaces, minimizing lock contention on shared memory
structures and enabling the tool to analyze enterprise-scale,
multi-million-line codebases in seconds.

\section{Inter-Language Dependency Detection}

Currently, while the framework successfully aggregates graphs
extracted from multi-language repositories into a unified
representation, the resolution of dependencies remains confined
within the boundaries of each individual language.

In modern distributed or highly optimized systems, cross-language
interactions are ubiquitous. A compelling future extension would
involve detecting and formally resolving \textbf{inter-language
dependencies}. For example, mapping Java Native Interface (JNI) calls
from a Java microservice to the underlying C/C++ native libraries, or
resolving Foreign Function Interface (FFI) bindings between Python
and Rust. Extending the resolution engine to bridge these language
boundaries would allow the extraction of a truly global architectural
dependency graph, offering unprecedented visibility into complex
polyglot architectures.

\section{Support for Alternative Programming Paradigms}

The empirical validation and heuristic design of the current
framework have primarily targeted procedural and object-oriented
programming paradigms (C-like languages, Java, Python, Rust). These
languages share foundational architectural concepts such as modules,
classes, structs, and explicitly invoked functions.

Although the Intermediate Representation ($\mathcal{D}_{IR}$) was
designed mathematically to be as agnostic as possible, extending the
framework to support entirely different programming paradigms
represents a significant future challenge. Future research should
evaluate and expand the model's capability to capture architectural
dependencies in \textbf{Functional} languages (e.g., Lisp, Haskell,
Clojure), which heavily rely on higher-order functions and closures,
or in \textbf{Logic} programming languages (e.g., Prolog).
Incorporating these paradigms will likely require introducing new
heuristic classifiers in the extraction phase and refining the query
algebra to accommodate exotic scoping rules and
pattern-matching-based execution semantics.
